\subsection{Exact evaluation of active score prefixes}
\label{sec:saturation}
At a fixed width, let $A$ count admitted original-positive positions and let
$B$ denote their admitted non-original buffer mass. Then $T=A+B$ and
$P=P_0+B/2$. We call a complete score threshold \emph{saturated} when
\begin{equation}
 A+B/2\ge P_0\qquad\hbox{and}\qquad H=K_w.
 \label{eq:saturation-condition}
\end{equation}
The condition includes the complete global numerical tie, even when only some
positions in that tie belong to the width's active support.

\begin{lemma}[Permanent saturation]
After a threshold satisfying~\eqref{eq:saturation-condition}, every subsequent
recall coordinate equals one. The remaining right-step PR contribution is
zero, and the remaining trapezoidal ROC contribution, including the explicit
terminal $(1,1)$ transition, is $1-F$, where $F$ is the current FPR.
\end{lemma}
\begin{proof}
Threshold relaxation cannot decrease $A$, $B$, or the number of hit expanded
components. Thus both conditions persist. The first is equivalent to
$(A+B)/(P_0+B/2)\ge1$; the second makes the existence factor one. Every later
PR recall increment vanishes. Every later ROC trapezoid has height one, so
its signed horizontal increments telescope from the current $F$ to the
terminal abscissa one. The proof does not assume a fixed negative denominator.
\end{proof}

Stopping at recall one is a standard flat-curve observation, not a claimed
new integration principle. In particular, admitting the last original
positive gives $A=P_0$ and hits every expanded component, so it is already a
valid cutoff. The VUS-specific condition can stop earlier: buffer mass can
compensate for original positives that have not yet been admitted. Both
last-original-positive stopping and disabled stopping are explicit controls.

A cutoff added after constructing every width's full support histogram would
still pay $\Theta(Q)$ per width. To avoid that work, retain support positions
in descending score order, and activate position $i$ at radius $d_1(i)$;
original positives have activation radius zero. The score-ordered list is
represented by integer next pointers. Bucket-sort activations by their
integer radius. Starting from the full list, delete positions in reverse
activation order and record the two neighbors at each deletion. Reversing
these deletions inserts every position in constant time. Tied activation
radii use the same fixed order in both directions. Consequently traversing
the current list visits only positions with nonzero weight at that width.
A complete retained event still uses its precomputed global $c^>,c^\ge$;
points omitted from the current list may belong to that same global tie.

For geometry, preprocessing also records one nearest original interval for
each support position. At radius $r$, neighboring original intervals with
endpoints $b_k,a_{k+1}$ merge precisely when
$2r\ge a_{k+1}-b_k$. Bucket these edges by
$\lceil(a_{k+1}-b_k)/2\rceil$ and maintain their components with union-find.
Mere adjacency of expanded intervals does not merge them. An active point
lies in its nearest interval's expansion. Any other expansion containing the
point overlaps that expansion and belongs to the same component; hence the
nearest interval's union-find root identifies the point's component. A
width-stamped root array counts first hits without clearing all components
at each width. Unions occur before a width's threshold scan, so roots are
fixed throughout that scan.

\begin{theorem}[Active-prefix bound]
Use the fixed-binary64 word-RAM and unit-cost metric-arithmetic assumptions
of Proposition~\ref{prop:radix}. Let $J_w$ be the number of active support
positions visited through the complete first saturated score tie at width
$w$, and let $K$ count original intervals. Both scalar integrals can be
computed in
\[
 O\!\left(n+L+\alpha(K)\sum_{w=0}^{L}J_w\right)
\]
time and $O(n+Q)$ auxiliary words, where $\alpha$ denotes the inverse
Ackermann factor for union-find. A diagnostic vector of per-width work
counts uses an additional $O(L)$ words. The values equal the target
integrals in exact metric arithmetic.
\end{theorem}
\begin{proof}
Fixed-width radix ranking, complete global tie counts, the nearest-interval
scan, and extraction of the ordered support take $O(n+Q+K)$ time.
All activation distances and finite merge radii are at most $n-1$, so
bucket arrays are capped at $\min(\lfloor L/2\rfloor,n-1)+1$, even if $L$
is much larger than $n$. Counting-sort schedules and reverse deletions take
$O(n+Q+K)$ time and space. Each position is activated once, and there are
at most $K-1$ unions. Union by size with path compression takes amortized
$O(\alpha(K))$ per operation. Every visited position performs one find;
event arithmetic and first-hit accounting take constant additional time.
No width requires an additional $Q$- or $K$-sized scan. Streaming the means
requires $L+1$ width iterations. Since at width zero saturation requires all
original positives, $J_0=P_0\ge K$, and the union cost is absorbed by
$\alpha(K)\sum_w J_w$. Active-event elimination uses the same negative-run
rectangles as before, and the lemma supplies the exact omitted suffix.
All scratch state is input-, support-, or interval-sized.
\end{proof}

The bound is adaptive, not a worst-case improvement for every score vector:
$J_w$ may equal the full active support size. Union-find and linked-list
access also have costs that a work-count reduction does not measure.
Nevertheless, a separating family exists even against the same active-list
algorithm stopped at the last original positive. For $L\ge2$, take
$n=4L+3$, with singleton originals at zero-based positions $L+1$ and $3L+1$.
Rank first the first original, then its two immediate neighbors, then one
immediate neighbor of the second original. Rank the second original last;
put all remaining scores strictly between the fourth and last scores.
At every width $w\ge2$, those three neighbors have weight at least
$1/\sqrt2$, so the first four active positions satisfy
\[
 A+B/2\ge 1+\frac{3}{2\sqrt2}>2=P_0
\]
and hit both expanded components. Thus $J_w=4$, while at widths zero and
one $J_w=2$. The last-original-positive control must traverse
$Q_w=4\lfloor w/2\rfloor+2$ active positions at every width, giving
$\Theta(L^2)$ traversal rather than $\Theta(L)$. With $K=2$ and $n=\Theta(L)$,
the full adaptive algorithm takes $O(n+L)$ on this family. This is a
structural operation-count separation, not an empirical runtime claim.

The implementation uses finite-precision arithmetic and can differ from a
full scan through summation order and suffix telescoping. The theorem is an
exact-arithmetic identity; measured numerical parity, floating-point limits,
and runtime comparisons must be reported separately.
